% Recursive train/generate/allocate/evaluate pipeline.
% Result-independent: this figure describes the experimental apparatus, not any
% outcome, so it is safe to include before results exist.
% Requires \usepackage{tikz} and \usetikzlibrary{arrows.meta,positioning} in main.tex.
\begin{tikzpicture}[
  node distance=9mm and 12mm,
  every node/.style={font=\small},
  stage/.style={
    draw, rounded corners=2pt, align=center,
    minimum height=8mm, minimum width=22mm, inner sep=4pt
  },
  frozen/.style={stage, fill=black!5},
  artifact/.style={stage, densely dashed},
  flow/.style={-{Stealth[length=2mm]}, thick},
  feedback/.style={-{Stealth[length=2mm]}, thick, densely dotted},
]

  \node[frozen] (manifests) {Frozen data\\manifests};
  \node[stage, right=of manifests] (train) {Train\\generation $g$};
  \node[stage, right=of train] (generate) {Generate\\synthetic corpus};
  \node[stage, right=of generate] (allocate) {Allocate human\\rescue $b_{g,m}$};

  \node[frozen, below=of manifests] (budget) {Lifetime budget $B$\\$\sum_{g,m} b_{g,m} = B$};
  \node[stage, below=of generate] (evaluate) {Held-out\\evaluation};
  \node[artifact, below=of evaluate] (chain) {Immutable chain\\result + manifest};
  \node[artifact, right=of chain] (analysis) {Chain-level\\statistics};

  \draw[flow] (manifests) -- (train);
  \draw[flow] (train) -- (generate);
  \draw[flow] (generate) -- (allocate);

  % The recursion: this generation's allocation feeds the next generation's training.
  \draw[feedback] (allocate.north) -- ++(0,6mm) -| node[pos=0.25, above] {generation $g{+}1$} (train.north);

  \draw[flow] (budget) -- (allocate.south -| budget.east) ;
  \draw[flow] (budget.east) -- (evaluate.west);
  \draw[flow] (generate.south) -- (evaluate.north);
  \draw[flow] (evaluate) -- (chain);
  \draw[flow] (chain) -- (analysis);

\end{tikzpicture}
